\chapter{Experimental Methodology}
\label{chap:experiments}

\section{Dataset}

The experiments use SST-2, the binary sentiment classification task derived from the Stanford
Sentiment Treebank \citep{socher2013recursive} and distributed as part of GLUE
\citep{wang2018glue}. Each input is a sentence and each label indicates positive or negative
sentiment. The task is suitable for controlled low-resource experiments because it has a standard
validation set and a simple accuracy metric, while still requiring lexical and compositional
judgment.

Training subsets contain 100, 500, 1,000, 5,000, or 10,000 examples. The full validation split is
used for evaluation. Because the subset is reshuffled by seed, reported standard deviations include
variation due to both initialization and sampled examples.

\section{Compared Methods}

\begin{description}[style=nextline]
  \item[Full fine-tuning] All model parameters are updated.
  \item[LoRA] Rank 8 adapters are attached to query and value projections.
  \item[AdaLoRA] Target rank 8 and initial rank 12 are used on query and value projections.
  \item[Star-LoRA] Rank, initial rank, scaling, dropout, and target modules are selected from the
  dataset profile, followed by PEFT AdaLoRA training.
  \item[Layer-wise heuristic] The Star-LoRA rank policy is retained, but additional modules are
  concentrated in upper layers.
  \item[Claude layer-wise] A Claude-generated layer map is used as an exploratory alternative.
\end{description}

\section{Primary Experiment}

The primary comparison is the archive \texttt{outputs/1e-4 various dt}. It contains full
fine-tuning, LoRA, AdaLoRA, and Star-LoRA for three seeds (13, 21, and 42), five sample budgets,
three epochs, batch size 8, float32 model weights, and learning rate $10^{-4}$ according to the
experiment directory and run configuration. Since old JSON artifacts do not independently
serialize the learning rate, the archive is kept intact rather than pooled with other sweeps.

\section{Secondary Five-Seed Sweep}

The archive \texttt{outputs/no change lr, various seeds} contains the same four methods over five
seeds. The current training script's default learning rate is $2\times10^{-5}$, and the directory
name indicates the unchanged or default learning-rate condition. Because the archived metadata
does not store learning rate, the sweep is presented as a separate robustness check and not as a
controlled single-factor learning-rate ablation.

\section{Layer-Wise and LLM Explorations}

The \texttt{rule\_base} archive contains the global rule-based method for five seeds up to 5,000
examples. The \texttt{pilot} archive contains heuristic layer-wise runs and directories labeled
Gemini. Since the Gemini requests failed and used heuristic fallback, only the heuristic rows are
interpreted. The \texttt{claude} archive contains genuine Claude-generated layouts at 5,000
examples for five seeds.

\section{Metrics}

Accuracy is
\begin{equation}
  \mathrm{Accuracy} = \frac{1}{M}\sum_{i=1}^{M}
  \mathbf{1}[\hat{y}_i=y_i].
\end{equation}
Macro F1 averages class-wise F1 and is useful when predictions collapse toward one class.
Runtime, samples per second, and steps per second measure efficiency. Mean and sample standard
deviation are reported over seeds. Selected rank shows whether the dataset-aware rule changes
capacity as intended.

\section{Threats to Experimental Validity}

The study uses one model family, one model size, and one classification dataset. Hyperparameters
are not tuned independently for every method. Star-LoRA changes rank, scaling, dropout, schedule,
and target modules together, so the main comparison does not isolate which component causes an
improvement. Runtime is measured on the available system but hardware metadata is absent.
Finally, three seeds in the primary sweep provide only a limited estimate of variance.

These limitations determine the strength of the claims. The experiments can demonstrate prototype
behavior and motivate further research; they cannot establish universal superiority.
